\chapter{数学基礎補足}
\label{app:math}

この付録では、本文で使用した数学的概念を補足します。微積分と基礎線形代数を知るが行列微分などに慣れていない読者のための参考資料だ。

\section{ベクトルと行列}

\keyterm{ベクトル}は数字の1次元配列です。$\mathbf{x} = [x\_1, x\_2, \dots, x\_n]^T$と表記。\keyterm{行列}は 2 次元配列。$A \in \mathbb{R}^{m \times n}$は$m$行$n$列。

\begin{itemize}
\item\textbf{内積（dot product)}：$\mathbf{a} \cdot \mathbf{b} = \sum\_i a\_i b\_i$。スカラー戻り。
\item\textbf{行列積}:$C = AB$where$C\_{ij} = \sum\_k A\_{ik} B\_{kj}$.
\item\textbf{転置（transpose)}:$A^T$は行と列を変更します。
\item\textbf{norm}：$\|\mathbf{x}\|\_2 = \sqrt{\sum\_i x\_i^2}$。ベクトルの長さ。
\end{itemize}

\section{Softmax}

\keyterm{Softmax}は、任意の実数ベクトルを確率分布に変換する関数です。

\begin{formula}[Softmax]
\[
\text{softmax}(x\_i) = \frac{e^{x\_i}}{\sum\_j e^{x\_j}}
\]
出力のすべての要素は0と1の間であり、合計は1です。アテンション重み、分類モデルの出力などに使用。
\end{formula}

\section{Cross-Entropy}

\keyterm{クロスエントロピー（cross-entropy)}は2つの確率分布$p, q$の差を測定します。

\begin{formula}[Cross-Entropy]
\[
H(p, q) = -\sum\_i p\_i \log q\_i
\]
分類の問題では、$p$は正解の1つのホットベクトル、$q$はモデル予測softmax出力です。単一の正解$y$に対して$H = -\log q\_y$に単純化。
\end{formula}

\section{連鎖法則と Backpropagation}

\keyterm{連鎖法則（chain rule)}は合成関数の微分法です。

\[
\frac{d}{dx} f(g(x)) = f'(g(x)) \cdot g'(x)
\]

ニューラルネットワークは合成関数の連続であるため、連鎖法則を使用して入力から出力までの微分を計算できます。これが\keyterm{backpropagation}の数学的基盤です。

PyTorchはこのプロセスを自動化するので、私たちが直接微分式を導く必要はありません。しかし、「どのパラメータが損失にどのように影響するか」を理解するには、連鎖法則の概念を知る必要があります。

\section{行列微分}

損失$\mathcal{L}$が行列$W$に依存する場合、$\partial \mathcal{L} / \partial W$は$W$と同じ形状の行列であり、各要素は$\partial \mathcal{L} / \partial W\_{ij}$です。

\keyterm{Jacobian}はベクトル関数の微分を表す行列です。$\mathbf{y} = f(\mathbf{x})$のとき$J\_{ij} = \partial y\_i / \partial x\_j$。

詳細については、「The Matrix Cookbook」（Petersen \& Pedersen）を参照してください。

\section{正規分布}

\keyterm{正規分布（normal distribution)}$\mathcal{N}(\mu, \sigma^2)$の確率密度関数:

\[
p(x) = \frac{1}{\sqrt{2\pi\sigma^2}} \exp\left(-\frac{(x-\mu)^2}{2\sigma^2}\right)
\]

重みの初期化によく使用されます。$\mu=0, \sigma = 1/\sqrt{d}$を書くと、入力次元が大きくなっても出力分散が一定に保たれます（Xavier初期化の基本的なアイデア）。

\section{情報理論の基礎}

\keyterm{エントロピー（entropy)}は確率分布の「ランダム性」を測定します。

\[
H(p) = -\sum\_i p\_i \log p\_i
\]

\keyterm{KL 発散（KL divergence)}は2つの分布の違いです。

\[
D\_{KL}(p \| q) = \sum\_i p\_i \log \frac{p\_i}{q\_i}
\]

Cross-entropyは$H(p) + D\_{KL}(p \| q)$に分解されます。$H(p)$が固定なので、クロスエントロピーを最小化することは$D\_{KL}$を最小化するのと同じです。

\section{最適化の基礎}

\keyterm{傾斜下降法（gradient descent)}：損失を減らす方向にパラメータを更新します。

\[
\theta \leftarrow \theta - \eta \nabla\_\theta \mathcal{L}
\]

$\eta$は学習率です。\keyterm{Adam}は、運動量と適応学習率を組み合わせた最適化アルゴリズムです。\keyterm{AdamW}はAdamにweight decayを組み合わせたバージョンで、LLM学習の標準です。

詳細な誘導は、Goodfellow et al。 「Deep Learning」4$\sim$第8章を参照してください。
